\documentclass[11pt]{article}
\usepackage[margin=1in]{geometry}
\usepackage{amsmath,amssymb,booktabs,hyperref}
\title{Replication Report: \emph{Microscopic origin of the spin-splitting in altermagnets} (Lee, Kim \& Kim, 2025)}
\author{Automated physics replication (Hermes subagent)}
\date{\today}
\begin{document}
\maketitle

\section{Headline claim tested}
The paper proposes a minimal two-dimensional square-lattice, four-band model
(basis $[A\!\uparrow,B\!\uparrow,A\!\downarrow,B\!\downarrow]^T$) whose spin
splitting arises from the interplay of an atomic exchange field $h_{\text{eff}}$
and an anisotropic third-neighbour hopping $\delta t$ generated by distortion of
non-magnetic ion cages. The testable headline: the momentum-dependent spin
splitting is nonzero \emph{only when both} $h_{\text{eff}}\neq0$ and
$\delta t\neq0$, with a $d$-wave nodal spin structure.

\section{Model (paper Eqs.\ 1--3)}
\begin{align}
H &= \epsilon_k\, I + t_{k,x}\,\tau_x + t_{k,z}\,\tau_z + h_{\text{eff}}\,(\tau_z\otimes\sigma_z),\\
\epsilon_k &= 2t_2(\cos k_x+\cos k_y) + 4t_3\cos k_x\cos k_y - \mu,\\
t_{k,x} &= 4t_1\cos(k_x/2)\cos(k_y/2),\qquad
t_{k,z} = -4\,\delta t\,\sin k_x\,\sin k_y,\\
\epsilon_{\alpha=\pm,\sigma} &= \epsilon_k \pm \sqrt{t_{k,x}^2 + (t_{k,z}+h_{\text{eff}}\sigma)^2}.
\end{align}
Parameters (energy unit $t_1$): $t_1{=}1,\ t_2{=}0.5,\ t_3{=}0.25,\ \delta t{=}0.2,\ h_{\text{eff}}{=}2.0,\ \mu{=}0$.

\section{Method}
We built the $4\times4$ Bloch Hamiltonian from scratch and diagonalized it
numerically (\texttt{numpy.linalg.eigh}), assigning spin labels via
$\langle\sigma_z\rangle$. Independent numeric bands were verified against the
closed-form Eq.~(3), then the spin splitting
$\Delta E(\mathbf{k})=|E_{+,\uparrow}-E_{+,\downarrow}|$ was scanned across
parameters and the Brillouin zone.

\section{Results}
\begin{table}[h]\centering
\begin{tabular}{lll}
\toprule
Check & Result & Status \\
\midrule
Numeric $4\times4$ vs analytic Eq.(3) & max err $1.8\times10^{-15}$ & PASS \\
Spin splitting at $(\pi/2,\pi/2)$ & $\Delta E = 1.109\,t_1$ & PASS \\
$\Delta E$ monotonic in $\delta t$ (Fig.\ 4d) & $0\to2.47$ over $\delta t\in[0,0.5]$ & PASS \\
$\Delta E$ monotonic in $h_{\text{eff}}$ (Fig.\ 4d) & $0\to1.43$ over $h\in[0,4]$ & PASS \\
Anisotropic split requires BOTH terms & only ``both\_on'' nonzero (1.174) & PASS \\
Nodal along $k_x{=}0$ and $k_y{=}0$ & max split $=0$ on both lines & PASS \\
$d$-wave sign change between quadrants & sign $t_{k,z}$: $-\!\to\!+$ & PASS \\
\bottomrule
\end{tabular}
\caption{Self-scored checks. 7/7 pass.}
\end{table}

The two-ingredient case study is decisive: with $\delta t{=}0.2,h_{\text{eff}}{=}0$
(distortion only) the maximum splitting over the BZ is $0$; with
$\delta t{=}0,h_{\text{eff}}{=}2$ (exchange only) the momentum-dependent
splitting is likewise $0$ (an isotropic AFM-like gap, no altermagnetic
anisotropy); only $\delta t{\neq}0$ \emph{and} $h_{\text{eff}}{\neq}0$ yields the
finite anisotropic splitting ($1.174\,t_1$). The splitting generator
$t_{k,z}\propto\sin k_x\sin k_y$ is a $d_{xy}$ form factor, giving nodes along
$k_x{=}0$/$k_y{=}0$ and a sign change between BZ quadrants --- the hallmark
$d$-wave nodal spin structure.

\section{Verdict}
\textbf{REPLICATED.} All seven quantitative checks of the minimal-model headline
claim reproduce the paper's stated behaviour. Coverage 8/10 (model core + all
qualitative claims replicated; DFT/MnF$_2$ quantitative fit and SOC/AHE
topology out of scope). Agreement 10/10 on the analytic band structure and the
two-ingredient necessity rule.

\section{Credits}
Kubo/Berry methodology reference: \texttt{gobel2024\_sd\_skyrmion\_kubo\_Lz\_kernel.py}
(shared-kernels-cache), cited for the anomalous-Hall follow-up; not required for
the SOC-free core here. Physics runner: \texttt{comfyui-env/bin/python}.
\end{document}
